\section{PDDL planning}
\label{sec:pddl}

\paragraph{Where PDDL sits.} PDDL is integrated as a \emph{member of the plan
library} (Section~\ref{sec:arch}), not as a replacement for the reactive loop.
The \texttt{go\_to} intention can be achieved either by the fast BFS plan
(\texttt{FollowPathGoTo}) or by a PDDL plan (\texttt{PddlGoTo}) generated from
current beliefs; both yield the same kind of deterministic move sequence,
executed by the same \texttt{ActionExecutor}. This keeps PDDL meaningful --- it
genuinely plans the \emph{means} of an intention --- while task selection (which
parcel, when to deliver) stays with the BDI strategy. The same domain also
expresses a full single-parcel collect-and-deliver subproblem
(\texttt{PddlPickUpAndDeliver}), used for experiments and gated off by default.

\paragraph{Domain.} The STRIPS domain models a single agent moving on the
directed tile graph and, optionally, picking up and delivering one parcel. One
key encoding choice: movement is a set of \emph{directed} edge predicates
(\texttt{up}/\texttt{down}/\texttt{left}/\texttt{right}), emitted only where
entry is allowed, so one-way arrow tiles are captured for free --- a forbidden
direction simply has no edge, and no special action is needed. Pickup and
putdown extend the same domain to a delivery task:

\begin{lstlisting}
(:predicates (at ?t) (parcel ?p) (parcel-at ?p ?t)
             (carrying ?p) (delivery ?t) (delivered ?p)
             (up ?from ?to) (down ?from ?to) (left ?from ?to) (right ?from ?to))

(:action move-up    :parameters (?from ?to)
  :precondition (and (at ?from) (up ?from ?to))
  :effect (and (not (at ?from)) (at ?to)))            ; one per direction

(:action pickup     :parameters (?p ?t)
  :precondition (and (parcel ?p) (at ?t) (parcel-at ?p ?t))
  :effect (and (not (parcel-at ?p ?t)) (carrying ?p)))

(:action putdown    :parameters (?p ?t)
  :precondition (and (parcel ?p) (at ?t) (delivery ?t) (carrying ?p))
  :effect (and (not (carrying ?p)) (delivered ?p)))
\end{lstlisting}

\paragraph{Problem generation from beliefs.} A problem is built on demand from
the BeliefBase. A BFS from the agent's tile collects the reachable region; each
tile becomes an object \texttt{t\_x\_y}, each allowed step an edge predicate.
For a path problem the goal is \texttt{(at \emph{target})}; for a delivery
problem we add the parcel object, \texttt{(parcel-at p t)}, every reachable
\texttt{(delivery t)} tile, and the goal \texttt{(delivered p)} --- letting the
solver sequence pickup, routing and delivery in one symbolic plan. The region is
capped at \texttt{PDDL\_MAX\_TILES} ($1600$) and the target/parcel must lie
inside it, otherwise generation fails fast and BFS takes over.

\paragraph{When PDDL fits.} PDDL is appropriate for
\emph{deterministic, fully-believed} subproblems and a poor fit for the outer
loop, for two measured reasons. First, the world changes per frame (parcels
decay, agents move), so a long symbolic plan goes stale within seconds.
Second, the online solver (\texttt{@unitn-asa/pddl-client}, planning-as-a-service)
adds a fixed latency of roughly $3$\,s per call --- even for a one-step plan.
We therefore call the solver only off the critical path, behind three guards:
calls are time-limited (\texttt{PDDL\_TIMEOUT\_MS}, default $2.5$\,s, then fall
back to BFS); short routes skip PDDL entirely (\texttt{PDDL\_MIN\_PATH\_LENGTH},
default $10$, since BFS is essentially instant there); and once parcels are
being carried, PDDL is skipped because reward decay makes the solver delay
especially expensive during delivery. Any failure
(timeout, unreachable target, missing package) raises and the next plan in the
library --- BFS --- serves the same intention, so PDDL is strictly optional at
runtime.

\paragraph{Measured comparison.} Table~\ref{tab:pddl} reports illustrative
$60$\,s live runs on the same scenario, contrasting BFS with PDDL movement and
the full PDDL delivery task. A naive use of PDDL --- calling the solver for
every \texttt{go\_to} --- collapsed the score to $0$: the plans were correct but
each cost $\sim 3$\,s, so the agent spent the run waiting instead of acting
($16$ planner calls, $0$ deliveries). The selective gate above changed the
movement result from $0$ to $137$ while cutting planner calls from $16$ to $4$.
The full pickup-and-deliver task plan works (parcels are collected and
delivered by a single solved plan) but remains solver-bound, so BFS/BDI stays
the pragmatic default and PDDL is kept as a meaningful, bounded component and as
report evidence.

\begin{table}[ht]
\centering\small
\begin{tabular}{@{}lrrrr@{}}
\toprule
Setup & Score & Delivered & Picked & Planner calls \\
\midrule
BFS / BDI (default)        & 231 & 21 & 21 & 0 \\
Naive PDDL movement        &   0 &  0 & 12 & 16 \\
Selective PDDL movement    & 137 & 13 & 14 & 4 \\
Full PDDL delivery task    &  36 &  5 &  8 & 14 \\
\bottomrule
\end{tabular}
\caption{BFS vs.\ PDDL on one $60$\,s scenario (single representative runs).
Planner calls and failures come from the \texttt{plannerCalls} metric and the
\texttt{pddl\_plan}/\texttt{pddl\_failure} logs. See Section~\ref{sec:experiments}
for the broader strategy comparison.}
\label{tab:pddl}
\end{table}

\paragraph{Problem complexity.} The generated problems are deliberately small
and tractable: objects are bounded by the reachable region ($\leq$\,\texttt{PDDL\_MAX\_TILES}),
edges are $O(\mathrm{tiles})$ since each tile has at most four neighbours, and the
path subproblem is shortest-path-equivalent (the same answer BFS computes in
$O(V{+}E)$). The cost is not search difficulty but the network round-trip to the
online solver; the delivery problem only adds one parcel and the delivery-tile
set, so it stays within the same complexity class while letting the planner
sequence pickup, routing and putdown together. Mission constraints
(deliver-exactly-$N$, value thresholds) are intentionally \emph{not} pushed into
PDDL: they are enforced by the BDI plans (Section~\ref{sec:bdi}), which keeps the
symbolic model simple and avoids regenerating problems as missions change. The
PDDL delivery plan stays \emph{mission-safe} by deferring to those BDI plans
whenever a mission constraint is active, rather than trying to satisfy every
mission type itself; as a guard it also reuses their compliant-subset putdown
selection, so it can never drop a non-compliant batch. It is therefore disabled
by default for latency, not for correctness.

\paragraph{Scope and future work.} Crate-pushing and dynamic-obstacle variants
(the Sokoban-style maps) are out of scope: the agent assumes a static known map
with partial sensing of parcels and opponents (Section~\ref{sec:bdi}). A natural
extension is a richer multi-parcel delivery domain with explicit capacity and
mission constraints.
